\section{Introduction}\label{sec:intro}

Reading text is one of the most frequent reasons blind people photograph their surroundings. The first crowd-powered visual question answering service, VizWiz, already found that a large share of questions asked what a label, sign or screen said~\cite{bigham2010vizwiz}, and later analyses of tens of thousands of such questions confirmed that reading and product identification dominate everyday visual needs~\cite{brady2013visual,gurari2018vizwiz}. Dedicated text readers such as Seeing AI and OrCam serve this need~\cite{granquist2021evaluation,khan2020insight}, and image-processing methods can adapt other visual content, such as illustrations, for low-vision readers~\cite{zatserkovnyi2024adapting}. Since 2023, general-purpose vision-language models (VLMs) have increasingly taken over the task. Scene-description apps built on these models let a blind user point a phone at a medicine box or a utility bill and receive a fluent spoken description~\cite{gonzalezpenuela2024investigating,xie2025beyond,adnin2024king}.

Fluency is the problem. A conventional OCR engine that cannot read a word usually produces visible garbage, and a sighted user would notice it. A VLM that cannot read a word tends to produce a \emph{plausible} word instead: a common brand name in place of an unusual one, a complete street name where the sign is cut off by the edge of the frame, or a date that is one digit off. For a blind user the description is often the only access to the text. Blind users are known to place considerable trust in machine-generated descriptions even when they are wrong~\cite{macleod2017understanding}, and they have few independent ways to check them~\cite{alharbi2024misfitting,hong2024understanding}. A hallucinated expiry date, dosage or amount due is therefore not a cosmetic error but an accessibility failure with practical consequences.

This risk is plausibly greater for Cyrillic text, and for Ukrainian in particular. Text-centric benchmarks report clearly lower VLM accuracy for non-Latin scripts~\cite{liu2024ocrbench,tang2025mtvqa}. Ukrainian adds letters that do not exist in Russian (\foreignlanguage{ukrainian}{і, ї, є, ґ}) and an apostrophe inside words, so a model with a strong Russian or generic Cyrillic prior can silently ``correct'' them. Ukrainian everyday text is also dense with proper names, such as brands, surnames, street names and invented shop names, which a language prior cannot predict.

This paper asks two questions:
\begin{description}
\item[RQ1] When a VLM describes a photograph of Cyrillic signage, packaging or a document for a blind user, how often does it assert text that is not printed there, and what drives these text-reading hallucinations?
\item[RQ2] Can such hallucinations be detected automatically, without ground truth, using signals that an assistive application could compute at run time?
\end{description}

Answering RQ1 requires knowing exactly what is printed and what is actually visible in the photograph. We address it in two studies. Study~1 is a controlled factorial experiment on \NImages{} synthetic, photograph-like images of Ukrainian signs, labels and documents. Every character and its visibility are known, and the design crosses image quality, framing (cropping by the frame edge, finger or glare occlusion) and the type of the distinctive name: a real word, an invented word, or a \emph{near-word} one letter away from a real word. Study~2 repeats the measurements on \RealNImages{} images of real Ukrainian text. These are \NPack{} packaging photographs from Open Food Facts and \NDoc{} official and public documents, such as Pension Fund and city-council forms, Ministry of Health medical forms, utility bills, bank forms, medicine leaflets and notices. Each has a ground truth that is independent of the evaluated models, and each is degraded and framed in the same way as in Study~1. A paired near-word probe edits one word on real document pages.

We evaluated four Claude models of one vendor (Haiku~5.5, Sonnet~5.5, Opus~5.5 and Fable~5.1) with a blind-user description prompt and a constrained transcription prompt. Every word the models quote as printed text was labelled automatically against the ground truth. The labeller was checked by two independent automated raters. The study is a comparison within one model family; it makes no claim about vision-language models in general.

The contributions are:
\begin{enumerate}
\item a synthetic factorial benchmark and a real-image benchmark of Ukrainian packaging and documents, both with character-level visibility ground truth, released with code, source lists and response-free results;
\item an error taxonomy grounded in visibility, which separates misreadings and fabrications of visible text from \emph{completions} of text the camera did not capture, with prevalence per model and condition;
\item evidence, on synthetic and on real documents, that the models normalise unfamiliar printed words to familiar ones, and measurements of how image quality, framing and word category drive errors, with paired tests across models;
\item an evaluation of run-time detectors with operating points, costs and latencies, and an assessment of run-to-run stability.
\end{enumerate}

\section{Related work}\label{sec:related}

\subsection{AI visual assistance for blind and low-vision people}

Remote sighted assistance and crowd-sourced answers~\cite{bigham2010vizwiz} were followed by datasets of real photographs taken by blind people, which made their characteristics measurable: VizWiz images are frequently blurred, poorly framed or partly occluded, and many of them contain text~\cite{gurari2018vizwiz,gurari2020captioning}. What blind and low-vision people want from image descriptions depends on the context, but transcribing visible text is consistently valued~\cite{stangl2020person}. Earlier computer-vision aids for visually impaired people focused on detection and recognition pipelines~\cite{zatserkovnyi2021cv,khan2020insight}. Text readers built on such pipelines are now an established tool, although their accuracy degrades on real-world photographs~\cite{granquist2021evaluation}.

Large multimodal models changed the interaction from reading out recognised text to conversational description. Diary and interview studies of apps such as Be My AI report broad adoption and new use cases, alongside errors that users find hard to detect and strategies for verifying output, such as re-taking photographs or asking again~\cite{gonzalezpenuela2024investigating,xie2025beyond,adnin2024king}. Studies of how blind people handle recognition errors show that verification is laborious and often impossible without sighted help~\cite{alharbi2024misfitting,hong2024understanding}. Blind users also tend to trust captions that are inconsistent with the context~\cite{macleod2017understanding}. These findings motivate detection that works \emph{inside} the assistive pipeline rather than relying on the user.

\subsection{Hallucination in vision-language models}

Object hallucination in image captioning was quantified early with the CHAIR metric, which attributed it partly to language priors~\cite{rohrbach2018object}. For instruction-tuned VLMs it has been probed with polling questions in POPE~\cite{li2023evaluating} and with adversarial visual questions in HallusionBench~\cite{guan2024hallusionbench}. Recent surveys organise the causes into data, model and decoding factors, and highlight the dominance of the language prior when visual evidence is weak~\cite{bai2024hallucination,liu2024survey}. We use the same intuition but measure it on text, where the visual evidence can be degraded in a controlled way. Our earlier work on inclusive image captioning with large language models~\cite{zatserkovnyi2025llm} examined description quality for visually impaired users. The present study isolates the text-reading component of such descriptions.

\subsection{Reading text with vision-language models}

Scene-text benchmarks such as TextVQA~\cite{singh2019towards}, ST-VQA~\cite{biten2019scene} and TextCaps~\cite{sidorov2020textcaps} established that answering and captioning often require reading. OCRBench showed that large multimodal models rely on semantic plausibility: they read meaningful words far better than non-semantic character strings~\cite{liu2024ocrbench}. MTVQA found a substantial accuracy gap for non-Latin scripts~\cite{tang2025mtvqa}. On food packaging, a controlled comparison under realistic print degradations found that vision-language models extract structured label information better than classical OCR, but that a small model exhibits a hallucination mode with high completeness and low accuracy~\cite{zatserkovna2026packaging}. Text in an image also strongly steers model behaviour, as typographic attacks show~\cite{cheng2024unveiling}. Closest to our work, recent studies document OCR hallucination by multimodal models on degraded documents~\cite{he2025seeing} and fluent fabrication by VLM-based OCR on historical documents even where the character error rate is low~\cite{gardella2026low}. None of these studies controls what is visible in the image or targets the assistive setting, Cyrillic script or Ukrainian orthography. We use synthetic rendering, a standard way to obtain exact text ground truth~\cite{jaderberg2016reading,gupta2016synthetic}, to provide this control.

\subsection{Detecting hallucinations without ground truth}

Black-box detectors exploit the observation that hallucinated content is unstable across samples: SelfCheckGPT compares a response with stochastic re-samples~\cite{manakul2023selfcheckgpt}, and semantic entropy measures disagreement among sampled meanings~\cite{kuhn2023semantic,farquhar2024detecting}. Models can also be asked to state their confidence~\cite{lin2022teaching,kadavath2022language}. However, verbalised confidence tends to be overconfident~\cite{xiong2023can}, and hedging language does not faithfully reflect a model's intrinsic uncertainty~\cite{yona2024can}. For text reading, an independent OCR engine such as Tesseract~\cite{smith2007overview} offers a cross-check that does not share the VLM's language prior. Our study compares these families of signals on the same quoted words.
